A la izquierda, los dos escenarios que repetimos con Docker. Las diferencias llegan al 21 por ciento en Folly. No es ruido despreciable.

A la derecha, la comprobación ampliada sin Docker, en los seis escenarios. Dos cosas importantes. La primera: el recorte no se equivoca al azar, tiene un sesgo. Con contención da más casi siempre, y con un millón de pilas da menos. No sabemos por qué. La segunda: el orden casi no cambia.

De aquí sale la regla que usamos en toda la presentación: el recorte sirve para el orden, no para el decimal. Una ventaja del diez por ciento no demuestra nada; una de dos o tres veces, sí.
